% --- Archon physics formalization source begin ---
% archon:physics
% archon:covers IPhO2026Problems/problem_IPhO_2026_2_B_3.lean
% archon:source-report reports/ipho_2026_k3/problem_IPhO_2026_2_B_3.source.json
% archon:problem-id IPhO_2026_2
% archon:part-id B.3
% archon:previous-part IPhO_2026_2_B_1
% archon:previous-part-policy natural_language_prerequisite_only
% archon:previous-part IPhO_2026_2_B_2
% archon:previous-part-policy natural_language_prerequisite_only

\chapter{Physics problem IPhO\_2026\_2\_B\_3}
\label{ch:IPhO2026Problems_problem_IPhO_2026_2_B_3}

\paragraph{Problem source.}
A half-cylindrical mirror of radius R illuminates a fully absorbing cylindrical
container of radius a.  Their axes are parallel, and the container center lies
R/2 from the mirror center on the symmetry plane.  Uniform parallel sunlight
arrives along the optical axis.  Any ray absorbed by the container reflects at
most once.  Let theta\_max be the largest incidence angle on the mirror among
rays that strike the container, and let P\_0 be the power the cylinder would
receive without the mirror.  See Figure 2f.

Current subquestion:
For R = 1.0 m, find a such that P = 5*P\_0, and report it in cm.

\paragraph{Current subquestion.}
For R = 1.0 m, find a such that P = 5*P\_0, and report it in cm.

\paragraph{Recorded answer/context.}
cos(theta\_max) = 4/5 and a = 0.12 m = 12 cm.

\paragraph{Figure/image path.}
/root/proposal\_for\_physic/science-mango/ipho\_2026\_source/image/T2\_page-3.png

\paragraph{Reusable previous-part conclusions.}
\begin{itemize}
\item Source B.1. Question: Given a = alpha*sin(theta\_max) + beta*sin(2*theta\_max), determine alpha and beta in terms of R. Reusable conclusions: alpha = R and beta = -R/2. Policy: natural\_language\_prerequisite\_only; do\_not\_import\_Lean\_output
\item Source B.2. Question: Express the power ratio P/P\_0 in terms of theta\_max. Reusable conclusions: P/P\_0 = 1/(1 - cos(theta\_max)). Policy: natural\_language\_prerequisite\_only; do\_not\_import\_Lean\_output
\end{itemize}

\paragraph{Formalization target.}
create a compiling Lean file with sorry bodies at `IPhO2026Problems/problem\_IPhO\_2026\_2\_B\_3.lean`.
The Lean declarations must preserve the physical quantities, dimensions or dimensional roles, figure labels, governing-law hypotheses, and final relation expressed by this problem.
Use Mathlib/Physlib names found through LeanExplore where available. If a domain API is missing, introduce faithful local abstractions rather than scalar placeholder aliases.

\begin{theorem}[Physics formalization target]
\label{thm:physics:IPhO_2026_2_B_3:target}
\uses{thm:IPhO2026Problems_problem_IPhO_2026_2_B_3:container_diameter_for_quintuple_power}
The assigned autoformalize agent should translate this physics problem into Lean declarations in the covered file, with theorem and lemma proof bodies written as `by sorry`.
\end{theorem}
\begin{proof}
This is an autoformalization task, not a proof task. Produce faithful statements that can later be proved without weakening the source contract.
\end{proof}
% NOTE: PhysLean-coverage exemption (planner-recorded, iter-002): PhysLean has no geometric/reflection-optics module for this mirror part (see this file's physics-grounding log). Self-containment is kept with the `import Mathlib` baseline; no irrelevant Physlib import is added. This documents the resolved import policy (iter-001 review ruling) for the blueprint-doctor `missing-physlib-import` check.
% --- Archon physics formalization source end ---

% --- Archon named-quantities coverage (blueprint-writer 2-b-3-entries) ---
% NOTE: import policy (mirrors the reconciliation note above, iter-002 exemption):
% the covered file builds on the `import Mathlib` baseline alone; PhysLean has no
% geometric/reflection-optics module for this mirror part, so no Physlib import is
% added.  This NOTE records the resolved import policy for the chapter.
% NOTE (scan-invisible pins): the covered file opens with `noncomputable section`,
% which hides its defining declarations (`abbrev`/`def` companions) from the hybrid
% decl scan; the `Prop`-valued structures and the theorems are scanned.  All ten
% entries below name live ROOT-level declarations (no namespace prefix on disk).
\subsection*{Unit scale and symbolic constants}

\begin{definition}[Recorded incidence angle]
\label{def:IPhO2026Problems_problem_IPhO_2026_2_B_3:thetaMaxRecorded}
\lean{thetaMaxRecorded}
Symbolic abbreviation for the maximal incidence angle of Figure~2f,
$\theta_{\max}$.  It is defined as the acute angle of the 3-4-5 triangle
recorded in the official solution.  This is a pure abbreviation: the
physical content of the official answer (the cosine of this angle) is
PROVED on the conclusion side of
\cref{thm:IPhO2026Problems_problem_IPhO_2026_2_B_3:container_diameter_for_quintuple_power},
not defined here.
\end{definition}
\begin{proof}
Definition; no claim.
\end{proof}

\begin{definition}[Metre in centimetres]
\label{def:IPhO2026Problems_problem_IPhO_2026_2_B_3:metreInCentimetres}
\lean{metreInCentimetres}
Unit scale: one metre expressed in centimetres, the conversion factor used
to report the final container radius in cm as the subquestion requests.
\end{definition}
\begin{proof}
Definition; no claim.
\end{proof}

\subsection*{Cross-section geometry}

\begin{definition}[Cross-sectional plane]
\label{def:IPhO2026Problems_problem_IPhO_2026_2_B_3:CrossSectionPlane}
\lean{CrossSectionPlane}
The Euclidean plane $\mathbb{E}^2$ containing the transverse cross-section
of the mirror/container system of Figure~2f.  The physical system is
translationally invariant along the cylinder axes, so a single cross-section
captures all the geometry; lengths in this plane are identified with real
numbers via the norm.
\end{definition}
\begin{proof}
Definition; no claim.
\end{proof}

\begin{definition}[Solar cooker geometry]
\label{def:IPhO2026Problems_problem_IPhO_2026_2_B_3:SolarCookerGeometry}
\lean{SolarCookerGeometry}
\uses{def:IPhO2026Problems_problem_IPhO_2026_2_B_3:CrossSectionPlane}
The geometric configuration of Figure~2f: the mirror radius $R$, the
container radius $a$, the container centre, and the unit sunlight direction
along the optical axis, with the bundled figure readouts $0 < a < R$ and the
problem datum that the container centre lies $R/2$ from the mirror centre on
the symmetry plane.
\end{definition}
\begin{proof}
Definition; no claim.
\end{proof}

\subsection*{Governing laws and previous-part interfaces}

\begin{definition}[Half-cylindrical mirror physics]
\label{def:IPhO2026Problems_problem_IPhO_2026_2_B_3:HalfCylindricalMirrorPhysics}
\lean{HalfCylindricalMirrorPhysics}
\uses{def:IPhO2026Problems_problem_IPhO_2026_2_B_3:SolarCookerGeometry, def:IPhO2026Problems_problem_IPhO_2026_2_B_3:CrossSectionPlane}
The governing-law carrier at a reflection point $x$ of the mirror: $x$ lies
on the mirror surface ($\|x\| = R$) on the illuminated half, the outgoing
ray is the law-of-reflection image of the incoming direction in the radial
normal line at $x$ (specular reflection at the cylindrical surface), the
incidence angle at $x$ is measured from the normal drawn at the point of
incidence and is related to the geometry by the cosine law, and it is acute
(the single-reflection/uniform-irradiance regime: any ray absorbed by the
container reflects from the mirror at most once).
\end{definition}
\begin{proof}
Definition; no claim.
\end{proof}

\begin{definition}[Previous-part results (B.1 and B.2 interfaces)]
\label{def:IPhO2026Problems_problem_IPhO_2026_2_B_3:PreviousPartResults}
\lean{PreviousPartResults}
The named natural-language prerequisite interfaces of subquestions B.1 and
B.2 packaged over an incidence angle $\theta \in (0, \pi/2)$ with $P, P_0 > 0$:
the B.1 geometric relation
$a = R \sin\theta - (R/2)\sin(2\theta)$
(the determined coefficients $\alpha = R$, $\beta = -R/2$) and the B.2 power
ratio $P/P_0 = 1/(1 - \cos\theta)$.  Nothing in this structure fixes the
value of $\theta$ or of $a$, so the B.3 conclusions are not smuggled in.
\end{definition}
\begin{proof}
Definition; no claim.
\end{proof}

\subsection*{Trigonometric certificates of the 3-4-5 triangle}

\begin{lemma}[The recorded angle is acute]
\label{lem:IPhO2026Problems_problem_IPhO_2026_2_B_3:thetaMaxRecorded_mem_Ioo}
\lean{thetaMaxRecorded_mem_Ioo}
\uses{def:IPhO2026Problems_problem_IPhO_2026_2_B_3:thetaMaxRecorded}
The recorded maximal incidence angle lies in the acute range:
$\theta_{\max} \in (0, \pi/2)$, as required of every ray on the
single-reflection branch that reaches the container.
\end{lemma}
\begin{proof}
The angle is an arccosine of a value lying strictly between $0$ and $1$, and
arccosine sends the interior of $[-1, 1]$ to $(0, \pi)$ in a strictly
decreasing way, so the value sits strictly between $\arccos 1 = 0$ and
$\arccos 0 = \pi/2$.
\end{proof}

\begin{lemma}[The sine of the recorded angle]
\label{lem:IPhO2026Problems_problem_IPhO_2026_2_B_3:sin_thetaMaxRecorded}
\lean{sin_thetaMaxRecorded}
\uses{def:IPhO2026Problems_problem_IPhO_2026_2_B_3:thetaMaxRecorded}
3-4-5 triangle certificate:
$\sin\theta_{\max} = 3/5$.
\end{lemma}
\begin{proof}
On the acute range of
\cref{lem:IPhO2026Problems_problem_IPhO_2026_2_B_3:thetaMaxRecorded_mem_Ioo},
$\sin\theta_{\max} = \sqrt{1 - \cos^2\theta_{\max}}$; the 3-4-5 ratio gives
$1 - (4/5)^2 = 9/25$, whose square root is $3/5$.
\end{proof}

\begin{lemma}[The sine of the doubled recorded angle]
\label{lem:IPhO2026Problems_problem_IPhO_2026_2_B_3:sin_two_mul_thetaMaxRecorded}
\lean{sin_two_mul_thetaMaxRecorded}
\uses{lem:IPhO2026Problems_problem_IPhO_2026_2_B_3:sin_thetaMaxRecorded, def:IPhO2026Problems_problem_IPhO_2026_2_B_3:thetaMaxRecorded}
Double-angle certificate: $\sin(2\theta_{\max}) = 24/25$.
\end{lemma}
\begin{proof}
Apply the double-angle identity $\sin(2\theta) = 2\sin\theta\cos\theta$ to the
3-4-5 ratio and
\cref{lem:IPhO2026Problems_problem_IPhO_2026_2_B_3:sin_thetaMaxRecorded}:
$2 \cdot (3/5) \cdot (4/5) = 24/25$.
\end{proof}

\subsection*{Target value theorem}

\begin{theorem}[Container radius for quintuple power]
\label{thm:IPhO2026Problems_problem_IPhO_2026_2_B_3:container_diameter_for_quintuple_power}
\lean{container_diameter_for_quintuple_power}
\uses{def:IPhO2026Problems_problem_IPhO_2026_2_B_3:SolarCookerGeometry, def:IPhO2026Problems_problem_IPhO_2026_2_B_3:HalfCylindricalMirrorPhysics, def:IPhO2026Problems_problem_IPhO_2026_2_B_3:PreviousPartResults, lem:IPhO2026Problems_problem_IPhO_2026_2_B_3:thetaMaxRecorded_mem_Ioo, lem:IPhO2026Problems_problem_IPhO_2026_2_B_3:sin_thetaMaxRecorded, lem:IPhO2026Problems_problem_IPhO_2026_2_B_3:sin_two_mul_thetaMaxRecorded, def:IPhO2026Problems_problem_IPhO_2026_2_B_3:metreInCentimetres, def:IPhO2026Problems_problem_IPhO_2026_2_B_3:thetaMaxRecorded}
\textbf{(Target, T2-B3.)}  For a solar-cooker geometry
(\cref{def:IPhO2026Problems_problem_IPhO_2026_2_B_3:SolarCookerGeometry})
governed by the mirror physics
(\cref{def:IPhO2026Problems_problem_IPhO_2026_2_B_3:HalfCylindricalMirrorPhysics})
at a reflection point reaching the container, with mirror radius
$R = 1\,\mathrm{m}$, container radius $a$, received power $P$ and no-mirror
power $P_0$ satisfying the previous-part relations B.1 and B.2
(\cref{def:IPhO2026Problems_problem_IPhO_2026_2_B_3:PreviousPartResults}),
if $P = 5\,P_0$ then: $\theta_{\max}$ equals the recorded angle
($\cos\theta_{\max} = 4/5$), the container radius is
$a = 0.12\,\mathrm{m}$, and reported in centimetres $a = 12\,\mathrm{cm}$.
The official answers $\cos\theta_{\max} = 4/5$ and $12\,\mathrm{cm}$ appear
here conclusion-side only; no hypothesis fixes them in advance.
\end{theorem}
\begin{proof}
Invert the B.2 ratio: $P/P_0 = 5 = 1/(1-\cos\theta_{\max})$ forces
$1-\cos\theta_{\max} = 1/5$, hence $\cos\theta_{\max} = 4/5$; since both
\cref{def:IPhO2026Problems_problem_IPhO_2026_2_B_3:PreviousPartResults} and
\cref{def:IPhO2026Problems_problem_IPhO_2026_2_B_3:HalfCylindricalMirrorPhysics}
confine $\theta_{\max}$ to the acute range $(0, \pi/2)$, on which cosine is
strictly monotone (so arccosine is its inverse there),
\cref{lem:IPhO2026Problems_problem_IPhO_2026_2_B_3:thetaMaxRecorded_mem_Ioo}
identifies $\theta_{\max}$ with the recorded angle.  Substitute into the B.1
relation $a = R\sin\theta_{\max} - (R/2)\sin(2\theta_{\max})$ at $R = 1$
with
\cref{lem:IPhO2026Problems_problem_IPhO_2026_2_B_3:sin_thetaMaxRecorded} and
\cref{lem:IPhO2026Problems_problem_IPhO_2026_2_B_3:sin_two_mul_thetaMaxRecorded}:
$a = 3/5 - (1/2)(24/25) = 3/25 = 0.12\,\mathrm{m}$, and the conversion scale
of
\cref{def:IPhO2026Problems_problem_IPhO_2026_2_B_3:metreInCentimetres} turns
this into $a \cdot 100 = 12$, i.e. $a = 12\,\mathrm{cm}$.
\end{proof}
